% !TeX root = ../main.tex
%% ------------------------------------------------------------------------- %%
\chapter{Implementação}
\label{cap:implementacao}

Os conceitos que sustentam a solução --- Arquitetura de Conjunto de Instruções (\textit{ISA}), Camada de Abstração de \textit{Hardware} (\textit{HAL}), ambiente de execução isolado (\textit{sandbox}) e \textit{Model-Based Design} --- foram apresentados no Capítulo~\ref{cap:fundamentacao}.
Este capítulo descreve como cada um deles se materializa na implementação: a arquitetura em camadas (Seção~\ref{sec:visao-geral-implementacao}), a plataforma-alvo (Seção~\ref{sec:plataforma-alvo}), a máquina virtual e o conjunto de instruções (Seção~\ref{sec:vm-isa}), o modelo de memória e o ambiente de execução isolado (Seção~\ref{sec:memoria-sandbox}), a camada de abstração de \textit{hardware} (Seção~\ref{sec:hal-implementacao}), o protocolo de comunicação (Seção~\ref{sec:protocolo}), a modelagem em \textit{Simulink}/\textit{Stateflow} (Seção~\ref{sec:stateflow-mbd}) e a estratégia de validação (Seção~\ref{sec:validacao-testes}).
A especificação completa do conjunto de instruções e os ensaios do bloco de controle PID constam, respectivamente, do \ref{ape:isa} e do \ref{ape:pid}.
Cada escolha de projeto é acompanhada de sua justificativa, reunidas no Quadro~\ref{qua:escolhas-projeto}.

\tikzset{
    bloco/.style={draw, rounded corners=2pt, align=center, font=\small, fill=gray!10, inner sep=4pt, minimum height=8mm},
    estado/.style={draw, rounded corners=6pt, align=center, font=\small, fill=gray!10, inner sep=4pt, minimum height=9mm, minimum width=2.2cm},
    seta/.style={-{Stealth[length=2.4mm]}, thick},
    rotulo/.style={font=\footnotesize, inner sep=1.5pt, fill=white}
}

%% ------------------------------------------------------------------------- %%
\section{Visão Geral da Arquitetura Implementada}
\label{sec:visao-geral-implementacao}

A solução é organizada em quatro módulos sobrepostos a uma fronteira de portabilidade, como mostra a Figura~\ref{fig:arquitetura-geral}.
O módulo de máquina virtual reúne o motor de execução e o conjunto de instruções; o módulo de mapas de memória guarda o estado de todos os elementos endereçáveis pelo programa (entradas, saídas, memórias auxiliares, temporizadores, contadores, blocos de controle e registradores de dados); o módulo de comunicação implementa o protocolo com a ferramenta de \textit{host}; e a camada de abstração de \textit{hardware} concentra todo o acesso aos periféricos do microcontrolador.
O programa do usuário é compilado para \textit{bytecode}, transferido ao dispositivo pelo protocolo serial e armazenado como dado somente leitura, de modo que a lógica de aplicação nunca acessa periféricos diretamente: ela é apenas percorrida pela máquina virtual, e esta, por sua vez, só interage com o \textit{hardware} através do contrato da camada de abstração.
Essa fronteira única entre o motor de execução e o microcontrolador é o que torna o motor de execução independente da plataforma, conforme discutido na Seção~\ref{sec:fundamentos-solucao}.

\begin{figure}[htb]
    \centering
    \begin{tikzpicture}[x=1cm,y=1cm]
        \node[bloco, fill=corApp, minimum width=10.4cm] (app) at (0,5.0) {Lógica de aplicação: \textit{bytecode} do usuário (dado somente leitura)};
        \node[bloco, fill=corVM, minimum width=3.0cm, minimum height=14mm, font=\footnotesize] (vm) at (-3.8,3.0) {Máquina virtual:\\motor de execução\\e instruções};
        \node[bloco, fill=corMem, minimum width=3.0cm, minimum height=14mm, font=\footnotesize] (map) at (0,3.0) {Mapas de memória:\\imagens de E/S\\e blocos funcionais};
        \node[bloco, fill=corCom, minimum width=3.0cm, minimum height=14mm, font=\footnotesize] (com) at (3.8,3.0) {Comunicação:\\enquadramento COBS\\e CRC16};
        \node[bloco, fill=corHAL, minimum width=10.4cm] (hal) at (0,1.1) {Camada de abstração de \textit{hardware} (contrato único com o microcontrolador)};
        \node[bloco, fill=corHW, minimum width=10.4cm] (mcu) at (0,-0.1) {Microcontrolador STM32F767ZI e periféricos};
        \node[bloco, fill=white, minimum width=2.4cm, minimum height=14mm] (host) at (8.6,3.0) {Ferramenta\\de \textit{host}};
        \draw[seta] (vm.north |- app.south) -- (vm.north);
        \draw[seta, <->] (vm) -- (map);
        \draw[seta, <-] (com.north |- app.south) -- node[rotulo, right] {carga} (com.north);
        \draw[seta] (vm.south) -- (vm.south |- hal.north);
        \draw[seta] (map.south) -- (map.south |- hal.north);
        \draw[seta] (com.south) -- (com.south |- hal.north);
        \draw[seta] (hal) -- (mcu);
        \draw[seta, <->] (com) -- node[rotulo, above] {serial} (host);
    \end{tikzpicture}
    \caption{Visão geral dos módulos do firmware e fronteira de portabilidade}
    \label{fig:arquitetura-geral}

    \fonte{Elaborado pelo autor (2026).}
\end{figure}

O laço principal do firmware repete indefinidamente quatro etapas: a leitura das entradas físicas com filtro de \textit{debounce} configurável, a execução de uma varredura completa do \textit{bytecode} pela máquina virtual, a aplicação da imagem de saídas ao \textit{hardware} e o atendimento da comunicação.
Esse ciclo de varredura (\textit{scan cycle}) reproduz, em software, o funcionamento cíclico de um CLP descrito na Seção~\ref{sec:tema} \cite{bolton:2015}.
O Quadro~\ref{qua:recursos} reúne a capacidade dos recursos disponíveis ao programa do usuário, todos definidos por constantes de compilação e, portanto, ajustáveis sem alteração do motor de execução.

\begin{quadro}[htb]
    \centering
    \begin{tabular}{|l|c|p{7.2cm}|}
        \hline
        Recurso & Quantidade & Observação \\
        \hline
        Entradas digitais & 10 & Com filtro de \textit{debounce} \\
        \hline
        Saídas digitais & 12 & --- \\
        \hline
        Entradas analógicas & 8 & Conversor de 12 \textit{bits} \\
        \hline
        Saídas PWM & 4 & Uma delas analógica pura (0 a 10 V) \\
        \hline
        Temporizadores & 32 & TON, TOF e TP; bases de 10 ms a 1 min \\
        \hline
        Contadores & 32 & CTU, CTD e CTUD; valor de 32 \textit{bits} com sinal \\
        \hline
        Blocos PID e RAMP & 16 + 16 & Instâncias independentes \\
        \hline
        Pilha lógica e pilha de blocos & 16 + 16 & Ramificações e blocos aninhados \\
        \hline
    \end{tabular}

    \caption{Recursos disponíveis ao programa do usuário}
    \label{qua:recursos}
    \fonte{Elaborado pelo autor (2026).}
\end{quadro}


%% ------------------------------------------------------------------------- %%
\section{Plataforma-Alvo}
\label{sec:plataforma-alvo}

A implementação tem como alvo o microcontrolador STM32F767ZI, de núcleo \textit{Arm Cortex-M7} com frequência de até 216 MHz e unidade de ponto flutuante \cite{st:ds11532}.
Quatro requisitos da solução justificam essa escolha.
O primeiro é a aritmética de ponto flutuante: 13 das instruções aritméticas e 6 das de comparação operam sobre \textit{F32}, e a unidade de ponto flutuante do núcleo as executa em \textit{hardware}, com tempo previsível e sem bibliotecas de emulação, o que preserva o determinismo discutido na Seção~\ref{sec:problema-literatura} \cite{liu:2000}.
O segundo é a memória: o dispositivo oferece até 2 MB de \textit{Flash} e 512 KB de SRAM, incluindo memórias acopladas ao núcleo (128 KB de dados e 16 KB de instruções) destinadas a dados e rotinas críticas de tempo real \cite{st:ds11532}, o que permite alocar os mapas de memória e o laço de varredura em regiões de acesso previsível.
O terceiro é a cobertura de periféricos: conversores analógico-digitais de 12 \textit{bits}, temporizadores, controladores de DMA, interfaces seriais e unidade de cálculo de CRC atendem a todos os grupos de funções do contrato da camada de abstração (Seção~\ref{sec:hal-implementacao}) \cite{st:ds11532,st:rm0410}.
O quarto é o ferramental: a placa NUCLEO-F767ZI integra o depurador/gravador ST-LINK, dispensando sonda externa e disponibilizando um canal serial para o \textit{host} \cite{st:um1974}, e a ferramenta \textit{STM32CubeMX} gera a inicialização dos periféricos a partir de uma descrição gráfica.
Como o desempenho do núcleo é superior ao necessário para o interpretador, espera-se que o tempo de varredura seja determinado pelo tamanho do programa do usuário e não pelo custo de interpretação, hipótese verificada pelos ensaios de desempenho descritos na Seção~\ref{sec:validacao-testes}.

A Figura~\ref{fig:esboco-placa} apresenta o esboço, em blocos, da conexão entre a placa e o circuito de entradas e saídas.
A atribuição individual de pinos é definida no projeto do circuito e não altera o contrato da camada de abstração.

\begin{figure}[htb]
    \centering
    \begin{tikzpicture}[x=1cm,y=1cm]
        \node[bloco, fill=corHW, minimum width=3.6cm, minimum height=34mm] (mcu) at (0,0) {\textbf{NUCLEO-F767ZI}\\STM32F767ZI\\[2pt]GPIO\quad ADC\\TIM (PWM)\\USART + DMA\\CRC};
        \node[bloco, fill=corCom, minimum width=3.0cm] (di) at (-5.6,1.1) {10 entradas\\digitais};
        \node[bloco, fill=corCom, minimum width=2.3cm] (dic) at (-5.6,-0.3) {Condicionamento\\de sinais};
        \node[bloco, fill=corCom, minimum width=3.0cm] (ai) at (-5.6,-1.7) {8 entradas\\analógicas};
        \node[bloco, fill=corMem, minimum width=3.0cm] (do) at (5.6,1.5) {12 saídas\\digitais};
        \node[bloco, fill=corMem, minimum width=3.0cm] (pwm) at (5.6,0.1) {3 saídas PWM};
        \node[bloco, fill=corMem, minimum width=3.0cm] (ao) at (5.6,-1.4) {1 saída analógica\\0 a 10 V};
        \node[bloco, fill=white, minimum width=3.6cm] (pc) at (0,-3.4) {Ferramenta de \textit{host}\\(PC)};
        \draw[seta] (di.east) -- ++(0.5,0) |- ([yshift=6mm]mcu.west);
        \draw[seta] (ai.east) -- ++(0.5,0) |- ([yshift=-8mm]mcu.west);
        \draw[seta] ([yshift=8mm]mcu.east) -- ++(0.5,0) |- (do.west);
        \draw[seta] (mcu.east) -- ++(0.5,0) |- (pwm.west);
        \draw[seta] ([yshift=-8mm]mcu.east) -- ++(0.5,0) |- (ao.west);
        \draw[seta, <->] (mcu.south) -- node[rotulo, right] {porta serial virtual (ST-LINK)} (pc.north);
    \end{tikzpicture}
    \caption{Esboço em blocos da conexão entre a placa-alvo e o circuito de entradas e saídas}
    \label{fig:esboco-placa}

    \fonte{Elaborado pelo autor (2026), com base em \cite{st:um1974}.}
\end{figure}


%% ------------------------------------------------------------------------- %%
\section{Máquina Virtual e Conjunto de Instruções}
\label{sec:vm-isa}

\subsection{Ciclo de execução}
\label{sub:ciclo-execucao}

O motor de execução implementa o ciclo \textit{fetch-decode-execute} apresentado na Seção~\ref{sec:fundamentos-solucao}.
A cada varredura, o contador de instruções é posicionado no início do \textit{bytecode}, a pilha lógica e o acumulador são zerados e a imagem de saídas retentivas é reaplicada; em seguida, cada instrução é lida da memória de programa (\textit{fetch}), tem o campo de opcode isolado (\textit{decode}) e é despachada para a rotina correspondente (\textit{execute}) até que a instrução de fim de varredura (\textit{EOS}) seja alcançada, como ilustra a Figura~\ref{fig:fetch-decode-execute}.
Um opcode fora do repertório válido interrompe a varredura e leva a máquina ao estado de erro (\textit{trap}), no qual a execução permanece bloqueada, com um pino de sinalização acionado, até que um novo \textit{bytecode} seja transferido.

\begin{figure}[htb]
    \centering
    \begin{tikzpicture}[x=1cm,y=1cm]
        \node[estado, fill=corHAL, minimum width=2.0cm] (ini) at (-6.6,0) {Início da\\varredura};
        \node[estado, fill=corApp, minimum width=2.0cm] (fetch) at (-3.3,0) {\textit{Fetch}};
        \node[estado, fill=corCom, minimum width=2.0cm] (decode) at (0,0) {\textit{Decode}};
        \node[estado, fill=corMem, minimum width=2.0cm] (exec) at (3.3,0) {\textit{Execute}};
        \node[estado, fill=corHAL, minimum width=2.0cm] (fim) at (6.6,0) {Fim da\\varredura};
        \node[estado, fill=corErr, minimum width=2.0cm] (trap) at (0,-2.4) {Erro (\textit{trap})};
        \draw[seta] (ini) -- (fetch);
        \draw[seta] (fetch) -- (decode);
        \draw[seta] (decode) -- node[rotulo, above=1pt, fill=none] {válido} (exec);
        \draw[seta] (exec) -- node[rotulo, above=1pt, fill=none] {\textit{EOS}} (fim);
        \draw[seta] (exec.north) -- ++(0,0.9) -| node[rotulo, pos=0.25, above] {demais instruções} (fetch.north);
        \draw[seta] (decode.south) -- node[rotulo, right] {inválido} (trap.north);
    \end{tikzpicture}
    \caption{Ciclo \textit{fetch-decode-execute} da máquina virtual}
    \label{fig:fetch-decode-execute}

    \fonte{Elaborado pelo autor (2026).}
\end{figure}

O despacho é feito por uma tabela de 128 posições, uma para cada valor possível do opcode de 7 \textit{bits}, cada uma apontando para a rotina da instrução.
Essa escolha, em vez de uma cadeia de comparações sequenciais, torna o tempo de despacho constante --- acesso a um vetor --- e independente de qual instrução é executada, característica desejável para a previsibilidade do ciclo de varredura \cite{liu:2000,patterson:2020}.
Antes da primeira varredura de cada programa carregado, o \textit{bytecode} é percorrido duas vezes: uma validação rejeita opcodes inválidos, e um pré-processamento extrai, uma única vez, as bases de tempo de temporizadores, blocos PID e blocos RAMP, retirando esse trabalho estático do laço de varredura.

\subsection{Formato e repertório de instruções}
\label{sub:formato-repertorio}

Toda instrução começa por uma palavra-base de 32 \textit{bits}, alinhada em palavra, que contém o opcode e o operando principal.
Essa primeira palavra segue sempre um de dois formatos, mostrados na Figura~\ref{fig:formato-instrucao}: o genérico, usado por instruções que endereçam um operando (tipo, índice e sub-tipo), e o dos blocos PID e RAMP, que identificam a instância e a base de tempo.
Por isso o estágio de decodificação é uniforme, pois lê sempre a mesma quantidade de \textit{bytes} para identificar a instrução \cite{patterson:2020}.
O \textit{bit} mais significativo do \textit{byte} de opcode é reservado à depuração e mascarado antes do despacho.

O tamanho total da instrução, entretanto, não é fixo: a palavra-base pode vir acompanhada de uma estrutura auxiliar, de dois tipos.
O primeiro tipo são palavras adicionais com constantes imediatas de 32 \textit{bits}, colocadas logo após a palavra-base.
Uma operação aritmética, de comparação ou de conversão com operando constante ocupa duas palavras, e o mesmo vale para um temporizador ou contador com valor programado constante.
O escalonamento (\texttt{SCALE}) é o caso extremo: carrega sempre os quatro limites de faixa e, se o valor de entrada também for constante, um quinto valor, ocupando cinco ou seis palavras (20 ou 24 \textit{bytes}).
O número de palavras adicionais é determinado pelo opcode e pelo campo de tipo da própria palavra-base, de modo que é conhecido já na decodificação e é sempre múltiplo de palavra, o que preserva o alinhamento e permite à validação percorrer o programa sem marcadores de fim de instrução.

O segundo tipo de estrutura auxiliar são instruções que preparam o estado de outra.
As instruções \texttt{PID}, \texttt{RAMP} e \texttt{PWM\_OUT} ocupam apenas a palavra-base, e os seus parâmetros (\textit{setpoint}, ganhos, limites, taxas, \textit{duty} e frequência) são gravados antes por instruções \texttt{STW}.
Da mesma forma, as entradas de temporizadores e contadores (habilitação, contagem, \textit{reset} e carga) são gravadas por instruções \texttt{ST} anteriores a \texttt{TIM} e \texttt{COUNT}.
Essa separação mantém a palavra-base pequena e o formato uniforme, ao custo de o comportamento de um bloco funcional depender de instruções que o precedem, o que orienta a composição dos programas de teste dos blocos (Seção~\ref{sec:validacao-testes}).
O \ref{ape:isa} traz o tamanho de cada instrução.

\begin{figure}[htb]
    \centering
    \begin{tikzpicture}[x=1cm,y=1cm, font=\footnotesize]
        \def\bw{0.44}
        \foreach \b/\x in {0/0, 1/3.52, 2/7.04, 3/10.56} {
            \node[font=\footnotesize] at ({\x+1.76},1.95) {byte \b};
        }
        \node[anchor=east] at (-0.2,1.0) {Genérico};
        \tikzset{campo/.style={draw, minimum height=8mm, font=\scriptsize, inner sep=0pt, align=center, anchor=south west}}
        \node[campo, fill=corErr, minimum width=0.44cm] at (0,0.6) {D};
        \node[campo, fill=corVM, minimum width=3.08cm] at (0.44,0.6) {opcode};
        \node[campo, fill=corApp, minimum width=1.76cm] at (3.52,0.6) {tipo};
        \node[campo, fill=corMem, minimum width=1.76cm] at (5.28,0.6) {índice alto};
        \node[campo, fill=corMem, minimum width=1.76cm] at (7.04,0.6) {índice baixo};
        \node[campo, fill=corCom, minimum width=1.76cm] at (8.80,0.6) {sub-tipo};
        \node[campo, fill=corHAL, minimum width=3.52cm] at (10.56,0.6) {reservado};
        \node[anchor=east] at (-0.2,-0.4) {PID e RAMP};
        \node[campo, fill=corErr, minimum width=0.44cm] at (0,-0.8) {D};
        \node[campo, fill=corVM, minimum width=3.08cm] at (0.44,-0.8) {opcode};
        \node[campo, fill=corMem, minimum width=1.76cm] at (3.52,-0.8) {instância};
        \node[campo, fill=corCom, minimum width=0.88cm] at (5.28,-0.8) {TB};
        \node[campo, fill=corHAL, minimum width=0.88cm] at (6.16,-0.8) {res.};
        \node[campo, fill=corHAL, minimum width=3.52cm] at (7.04,-0.8) {reservado};
        \node[campo, fill=corHAL, minimum width=3.52cm] at (10.56,-0.8) {reservado};
        \node[anchor=west, font=\scriptsize, align=left] at (0,-1.5) {D: \textit{bit} de depuração (1); opcode (7); tipo, índices, sub-tipo e instância (4 \textit{bits} cada);\\ TB: base de tempo (2); res.: reservado (2); reservado (8).};
    \end{tikzpicture}
    \caption{Formato da palavra-base de instrução (32 \textit{bits})}
    \label{fig:formato-instrucao}

    \fonte{Elaborado pelo autor (2026).}
\end{figure}

O repertório contém 76 instruções, com opcodes de \texttt{0x00} a \texttt{0x4B}, além da instrução de fim de varredura (\texttt{0x7F}), organizadas nos grupos do Quadro~\ref{qua:grupos-instrucoes}.
A lógica de contatos e bobinas segue o modelo de \textit{Instruction List} da norma IEC 61131-3 \cite{iec:2025}, com pilha lógica para ramificações e pilha de blocos para associações série e paralelo de trechos de \textit{rung}.
A especificação de cada opcode, dos tipos de operando e dos sub-tipos de cada bloco funcional consta do \ref{ape:isa}.

\begin{quadro}[htb]
    \centering
    \begin{tabular}{|p{6.1cm}|c|c|}
        \hline
        Grupo & Instruções & Opcodes \\
        \hline
        Controle de fluxo e pilhas & 7 & \texttt{0x00}, \texttt{0x01}, \texttt{0x12}--\texttt{0x16} \\
        \hline
        Contatos e lógica booleana (incluindo bordas) & 15 & \texttt{0x03}--\texttt{0x11} \\
        \hline
        Armazenamento e bobinas (\texttt{ST}, \texttt{OUT}, \texttt{SET}, \texttt{RESET}, \texttt{TOGGLE}) & 5 & \texttt{0x02}, \texttt{0x17}--\texttt{0x1A} \\
        \hline
        Temporizadores e contadores & 2 & \texttt{0x1B}--\texttt{0x1C} \\
        \hline
        Aritmética e movimentação (\textit{S32}/\textit{F32}) & 26 & \texttt{0x1D}--\texttt{0x36} \\
        \hline
        Comparação (\textit{S32}/\textit{F32}) & 12 & \texttt{0x37}--\texttt{0x42} \\
        \hline
        Conversão e escalonamento & 4 & \texttt{0x43}--\texttt{0x46} \\
        \hline
        Saídas analógica e PWM, escrita em bloco e blocos PID e RAMP & 5 & \texttt{0x47}--\texttt{0x4B} \\
        \hline
        Total & 76 & \texttt{0x00}--\texttt{0x4B} \\
        \hline
    \end{tabular}

    \caption{Grupos de instruções da máquina virtual}
    \label{qua:grupos-instrucoes}
    \fonte{Elaborado pelo autor (2026).}
\end{quadro}

\subsection{Blocos funcionais com estado}
\label{sub:blocos-funcionais}

Blocos funcionais mantêm estado entre varreduras e, por isso, não o armazenam no \textit{bytecode}, que é somente leitura, mas em mapas de memória dedicados (Seção~\ref{sec:memoria-sandbox}).
Os temporizadores implementam os três tipos da IEC 61131-3 --- \textit{TON} (retardo na energização), \textit{TOF} (retardo na desenergização) e \textit{TP} (pulso) --- com tempo programado constante ou lido de registrador e quatro bases de tempo (10 ms, 100 ms, 1 s e 1 min).
A contagem parte da borda de subida de \texttt{EN} no \textit{TON} e no \textit{TP} (este não retentável durante o pulso) e da borda de descida no \textit{TOF}, e é medida contra a base de tempo do sistema, e não pela contagem de varreduras, para que o intervalo não dependa do tamanho do programa.
Alterações do tempo programado durante a contagem são absorvidas sem reiniciá-la.
Os contadores implementam \textit{CTU}, \textit{CTD} e \textit{CTUD}, com entradas de contagem sensíveis a borda, entradas de \textit{reset} e de carga (\textit{load}) e valor programado constante ou de registrador.
O bloco \textit{RAMP} é um limitador de taxa de variação: sua saída persegue o valor alvo com taxas máximas independentes de subida e de descida, por unidade de tempo configurável; com \texttt{EN} desligado a saída é zero, e na borda de subida de \texttt{EN} ela assume imediatamente o valor inicial, sem rampa.

O bloco PID implementa o algoritmo posicional na forma paralela, com quatro decisões de projeto justificadas pelo comportamento em malha fechada \cite{astrom:2006}.
Primeiro, o intervalo de amostragem é fixo e configurável, e não medido: o bloco só recalcula quando esse período decorre, o que impede que a variação do tempo de varredura entre nos termos integral e derivativo.
Segundo, o termo derivativo é calculado sobre a variável de processo e não sobre o erro, evitando o pico de saída (\textit{derivative kick}) quando o \textit{setpoint} muda.
Terceiro, o anti-\textit{windup} é feito por integração condicional (\textit{clamping}): a integração é suspensa quando a saída já está saturada e o erro a empurra mais fundo na mesma direção, e retomada assim que o erro inverte de sinal.
Quarto, no modo manual o integrador acompanha a saída (\textit{back-tracking}), de modo que a passagem de manual para automático ocorre sem salto (\textit{bumpless transfer}).
A banda morta é aplicada ao erro antes dos termos proporcional e integral, e a saída é sempre limitada ao intervalo configurado, inclusive no modo manual.


%% ------------------------------------------------------------------------- %%
\section{Modelo de Memória e Ambiente de Execução Isolado}
\label{sec:memoria-sandbox}

Toda a memória do sistema é alocada estaticamente em tempo de compilação: não há alocação dinâmica nem coleta de lixo, o que elimina fragmentação e a variação de tempo que esses mecanismos introduziriam \cite{liu:2000}.
Cada tipo de elemento endereçável possui um mapa dedicado --- imagens de entradas, saídas e memórias auxiliares (\textit{bits}), registradores de dados de 32 \textit{bits}, temporizadores, contadores, instâncias PID e RAMP, e entradas e saídas analógicas ---, e o programa do usuário os referencia apenas por tipo e índice, nunca por ponteiro, como mostra a Figura~\ref{fig:mapa-memoria}.

\begin{figure}[htb]
    \centering
    \begin{tikzpicture}[x=1cm,y=1cm]
        \node[bloco, fill=corVM, minimum width=3.6cm, minimum height=60mm, font=\small] (vm) at (0,0) {\textbf{Máquina virtual}\\[6pt]\footnotesize acesso por tipo\\e índice, validado\\a cada operação};
        \node[bloco, fill=corApp, minimum width=6.8cm, minimum height=9mm] (bc) at (7.2,2.4) {\textit{Bytecode} (somente leitura)};
        \node[bloco, fill=corMem, minimum width=6.8cm, minimum height=9mm] (img) at (7.2,1.2) {Imagens de entradas, saídas e memórias (\textit{bits})};
        \node[bloco, fill=corHW, minimum width=6.8cm, minimum height=9mm] (dr) at (7.2,0) {Registradores de dados (32 \textit{bits})};
        \node[bloco, fill=corRos, minimum width=6.8cm, minimum height=9mm] (fb) at (7.2,-1.2) {Temporizadores, contadores, PID e RAMP};
        \node[bloco, fill=corCom, minimum width=6.8cm, minimum height=9mm] (an) at (7.2,-2.4) {Entradas analógicas e saídas PWM};
        \draw[seta, <-] (vm.east |- bc.west) -- (bc.west);
        \draw[seta, <->] (vm.east |- img.west) -- (img.west);
        \draw[seta, <->] (vm.east |- dr.west) -- (dr.west);
        \draw[seta, <->] (vm.east |- fb.west) -- (fb.west);
        \draw[seta, <->] (vm.east |- an.west) -- (an.west);
    \end{tikzpicture}
    \caption{Regiões de memória e forma de acesso pela máquina virtual}
    \label{fig:mapa-memoria}

    \fonte{Elaborado pelo autor (2026).}
\end{figure}

Essa segregação materializa o ambiente de execução isolado descrito na Seção~\ref{sec:fundamentos-solucao}.
O \textit{bytecode} é rejeitado se contiver opcode inválido, e as funções de acesso de cada mapa validam o índice a cada operação, e uma leitura fora da faixa retorna um valor neutro (falso ou zero) sem acessar memória, de modo que um erro de endereçamento na lógica do usuário não alcança memória fora das regiões alocadas nem o estado interno da máquina.
Diferentemente da compilação para código nativo, na qual um erro de lógica executa diretamente no processador, aqui o pior efeito possível é local à instrução defeituosa ou, no caso de opcode inválido, a parada controlada da execução com sinalização externa.


%% ------------------------------------------------------------------------- %%
\section{Camada de Abstração de Hardware}
\label{sec:hal-implementacao}

O contrato da camada de abstração expõe somente funções e tipos primitivos, organizados por grupo de periférico, e nenhum tipo do fabricante do microcontrolador aparece nessa interface: eles ficam confinados à implementação.
É essa separação entre contrato e implementação que faz da camada o único ponto do firmware dependente do microcontrolador, e que reduz a adaptação a outra plataforma à reimplementação dela.
O Quadro~\ref{qua:hal} associa cada grupo de funções ao periférico do STM32F767ZI que o implementa.

\begin{quadro}[htb]
    \centering
    \begin{tabular}{|p{5.6cm}|p{8.6cm}|}
        \hline
        Grupo de funções & Periférico e observações \\
        \hline
        Entradas e saídas digitais; sinalização de erro & \textit{GPIO} \\
        \hline
        Comunicação serial com o \textit{host} & \textit{USART} com DMA em transmissão e recepção, sobre a porta serial virtual do ST-LINK \\
        \hline
        Base de tempo (ms e $\mu$s) e temporizador de varredura & \textit{Timer} de propósito geral \\
        \hline
        Entradas analógicas (8 canais) & ADC de 12 \textit{bits} \\
        \hline
        Saídas PWM (4 canais) e saída analógica & \textit{Timer} em modo PWM; filtro externo para a saída de 0 a 10 V \\
        \hline
        Verificação de integridade (CRC16) & Unidade de CRC do microcontrolador ou rotina em \textit{software}, selecionada na inicialização \\
        \hline
    \end{tabular}

    \caption{Grupos de funções da camada de abstração e periféricos do STM32F767ZI}
    \label{qua:hal}
    \fonte{Elaborado pelo autor (2026), com base em \cite{st:ds11532,st:rm0410,st:um1974}.}
\end{quadro}

A implementação utiliza a biblioteca de \textit{HAL} fornecida pelo fabricante (\textit{STM32Cube HAL}), em vez do acesso direto a registradores (\textit{Low-Layer}), por dois motivos.
O código de inicialização gerado pela ferramenta \textit{STM32CubeMX} reduz o risco de erro de configuração de registradores, cuja depuração consome tempo desproporcional ao ganho de desempenho da alternativa de baixo nível nesta aplicação; e, como a camada de abstração do projeto já é uma fronteira própria, o eventual sobrecusto da biblioteca do fabricante fica contido nela e não se propaga ao motor de execução.
Interfaces de armazenamento externo e de USB não fazem parte do escopo desta monografia.


%% ------------------------------------------------------------------------- %%
\section{Protocolo de Comunicação}
\label{sec:protocolo}

O protocolo transfere o \textit{bytecode} do \textit{host} ao dispositivo e permite a depuração em tempo real (observação e forçamento de variáveis, pausa e retomada da execução).
Cada mensagem é um quadro composto pelo comando, um identificador de \textit{hardware}, os dados e uma verificação de integridade CRC16 de polinômio 0x8005 (forma refletida 0xA001), e é enviada com enquadramento \textit{Consistent Overhead Byte Stuffing} (\textit{COBS}).
O \textit{COBS} elimina o byte zero do conteúdo do quadro com sobrecarga limitada e conhecida, o que permite usar esse byte como delimitador sem ambiguidade \cite{cheshire:1999}; o CRC16 detecta corrupção em canais seriais ruidosos.
O identificador de dispositivo presente em cada quadro permite ao \textit{firmware} rejeitar comandos endereçados a um dispositivo diferente do esperado.
Os pacotes de dados têm no máximo 512 \textit{bytes}, dos quais até 256 de carga útil.

Os comandos organizam-se em três grupos: controle (confirmação positiva, confirmação negativa com código de erro e notificação de erro da máquina virtual), transferência (início, pacote de dados e fim) e depuração (início e fim, pausa e retomada, publicação e remoção de variáveis observadas, forçamento de valores e escrita direta em memória e registradores).
O dispositivo mantém uma máquina de estados com três modos --- espera, transferência e depuração ---, e em cada um deles rejeita, com um código de \textit{NACK} específico, qualquer comando inválido, de tamanho inconsistente, com CRC incorreto ou destinado a outro identificador de \textit{hardware}.
A Figura~\ref{fig:protocolo-sequencia} apresenta a sequência de uma transferência de \textit{bytecode}.
A autenticação do \textit{host} não faz parte do escopo desta monografia; a rejeição correta de quadros malformados é, portanto, a barreira de robustez do canal contra ruído e contra uso incorreto da ferramenta.

\begin{figure}[htb]
    \centering
    \begin{tikzpicture}[x=1cm,y=1cm]
        \node[bloco, fill=corApp, minimum width=3cm] (h) at (-4,0) {Ferramenta de \textit{host}};
        \node[bloco, fill=corVM, minimum width=3cm] (d) at (4,0) {Dispositivo};
        \draw[thick] (h.south) -- (-4,-6.4);
        \draw[thick] (d.south) -- (4,-6.4);
        \draw[seta] (-4,-1.0) -- node[rotulo, above] {\texttt{START\_TRANSFER}} (4,-1.0);
        \draw[seta, dashed] (4,-1.7) -- node[rotulo, above] {\texttt{ACK} (CRC e identificador válidos) ou \texttt{NACK}} (-4,-1.7);
        \draw[seta] (-4,-2.7) -- node[rotulo, above] {\texttt{DATA\_PACKET} $n$ (COBS + CRC16)} (4,-2.7);
        \draw[seta, dashed] (4,-3.4) -- node[rotulo, above] {\texttt{ACK} do pacote $n$ ou \texttt{NACK}} (-4,-3.4);
        \node[rotulo, fill=none] at (0,-4.05) {$\vdots$ repetido para cada pacote};
        \draw[seta] (-4,-4.8) -- node[rotulo, above] {\texttt{STOP\_TRANSFER}} (4,-4.8);
        \draw[seta, dashed] (4,-5.5) -- node[rotulo, above] {\texttt{ACK} ou \texttt{NACK}} (-4,-5.5);
    \end{tikzpicture}
    \caption{Sequência de uma transferência de \textit{bytecode}}
    \label{fig:protocolo-sequencia}

    \fonte{Elaborado pelo autor (2026).}
\end{figure}


%% ------------------------------------------------------------------------- %%
\section{Modelagem em Simulink/Stateflow segundo Model-Based Design}
\label{sec:stateflow-mbd}

A metodologia de \textit{Model-Based Design}, fundamentada na Seção~\ref{sec:fundamentos-solucao}, prescreve um modelo executável do comportamento esperado como referência para validar a implementação \cite{nicolescu:2009}.
Neste trabalho, os modelos de referência cumprem três papéis: servem de oráculo para os ensaios de lógica de contatos, temporização e contagem (Seção~\ref{sec:validacao-testes}), à maneira do modelo semântico executável empregado por Sadolewski e Trybus \cite{sadolewski:2024}; especificam a máquina de modos do bloco PID; e, em conjunto com modelos de planta, fecham a malha dos ensaios \textit{Hardware-in-the-Loop} do \ref{ape:pid}.
O procedimento de validação cruzada, análogo ao de Duggan et al. \cite{duggan:2022}, aplica o mesmo vetor de estímulos ao modelo simulado e à máquina virtual em execução, registra a imagem de saídas de cada varredura e compara os dois resultados bit a bit, com as transições de estado --- e não apenas o valor final --- incluídas na comparação.

A máquina de modos do bloco PID foi escolhida como primeiro artefato de modelagem por estar inteiramente especificada por estados e transições: quatro modos com prioridade estrita, avaliados nesta ordem a cada execução do bloco --- \textsc{Desabilitado}, \textsc{Reset}, \textsc{Manual} e \textsc{Automático} ---, como mostram a Figura~\ref{fig:stateflow-pid} e o Quadro~\ref{qua:estados-pid}.

\begin{figure}[htb]
    \centering
    \begin{tikzpicture}[x=1cm,y=1cm]
        \node[diamond, draw, aspect=2.2, align=center, font=\footnotesize, inner sep=1pt, fill=corHW] (en) at (0,0) {\texttt{EN}?};
        \node[diamond, draw, aspect=2.2, align=center, font=\footnotesize, inner sep=1pt, fill=corHW] (rs) at (4.6,0) {\texttt{RESET}?};
        \node[diamond, draw, aspect=2.2, align=center, font=\footnotesize, inner sep=1pt, fill=corHW] (mn) at (9.2,0) {\texttt{MANUAL}?};
        \node[estado, fill=corHAL] (des) at (0,-2.6) {\textsc{Desabilitado}};
        \node[estado, fill=corErr] (res) at (4.6,-2.6) {\textsc{Reset}};
        \node[estado, fill=corCom] (man) at (9.2,-2.6) {\textsc{Manual}};
        \node[estado, fill=corMem] (aut) at (12.6,0) {\textsc{Automático}};
        \draw[seta] (en) -- node[rotulo, above] {sim} (rs);
        \draw[seta] (rs) -- node[rotulo, above] {não} (mn);
        \draw[seta] (mn) -- node[rotulo, above] {não} (aut);
        \draw[seta] (en) -- node[rotulo, left] {não} (des);
        \draw[seta] (rs) -- node[rotulo, left] {sim} (res);
        \draw[seta] (mn) -- node[rotulo, left] {sim} (man);
    \end{tikzpicture}
    \caption{Prioridade dos modos do bloco PID (especificação do \textit{chart} \textit{Stateflow})}
    \label{fig:stateflow-pid}

    \fonte{Elaborado pelo autor (2026).}
\end{figure}

\begin{quadro}[htb]
    \centering
    \begin{tabular}{|l|p{4.6cm}|p{7.2cm}|}
        \hline
        Modo & Condição & Comportamento \\
        \hline
        \textsc{Desabilitado} & \texttt{EN} = falso & Bloco não calcula; \texttt{CV} assume o limite mínimo de saída; o integrador é preservado, para retomada próxima do ponto de parada. \\
        \hline
        \textsc{Reset} & \texttt{EN} verdadeiro e \texttt{RESET} verdadeiro & Integrador zerado; \texttt{CV} igual ao limite mínimo. \\
        \hline
        \textsc{Manual} & \texttt{EN} verdadeiro, \texttt{RESET} falso e \texttt{MANUAL} verdadeiro & \texttt{CV} segue o valor manual, limitado à faixa; integrador em \textit{back-tracking} para transição sem salto. \\
        \hline
        \textsc{Automático} & \texttt{EN} verdadeiro, \texttt{RESET} e \texttt{MANUAL} falsos & PID posicional com banda morta e anti-\textit{windup} por integração condicional. \\
        \hline
    \end{tabular}

    \caption{Especificação da máquina de modos do bloco PID}
    \label{qua:estados-pid}
    \fonte{Elaborado pelo autor (2026).}
\end{quadro}


%% ------------------------------------------------------------------------- %%
\section{Estratégia de Validação e Testes}
\label{sec:validacao-testes}

A validação é organizada em níveis de teste, escolhidos de acordo com o que cada um comprova isoladamente, conforme a prática da engenharia de \textit{software} descrita na norma ISO/IEC/IEEE 29119 \cite{iso:29119-1}.
Os \textbf{testes de unidade} isolam um grupo de instruções e comparam a imagem de saídas produzida com a de um oráculo.
Os \textbf{testes de integração} exercitam o protocolo de comunicação, o carregamento do \textit{bytecode} e a execução conjuntamente, inclusive com entradas malformadas.
Os \textbf{testes de sistema} fecham a malha de controle com uma planta simulada em tempo real (\textit{Hardware-in-the-Loop}, \textit{HIL}), técnica que permite validar o controlador real sem expor um processo físico a falhas \cite{isermann:1999}; sua aplicação ao bloco PID é apresentada no \ref{ape:pid}.
Os \textbf{testes de desempenho} medem o tempo de varredura e o consumo de memória, atendendo ao objetivo de análise quantitativa da solução.

A infraestrutura de teste é composta por quatro elementos.
Programas de teste são gerados por uma biblioteca de montagem de \textit{bytecode} em \textit{Python}, com uma função por instrução que codifica a palavra-base e as palavras auxiliares, o que dispensa o cálculo manual dos campos.
\textit{Scripts} no \textit{host} transferem o \textit{bytecode}, forçam valores de entrada e observam a imagem de saídas por meio dos comandos de depuração do protocolo.
Modelos em \textit{Simulink}/\textit{Stateflow} e cálculos em \textit{MATLAB} fornecem o resultado esperado (oráculo), e diagramas de temporização da própria norma IEC 61131-3 servem de referência para temporizadores, contadores e bordas \cite{iec:2025}.
Por fim, um instrumento de aquisição analógica (\textit{Analog Discovery 3}) mede as saídas analógicas e fecha a malha nos ensaios \textit{HIL}.
A Figura~\ref{fig:fluxo-validacao} resume o fluxo de um caso de teste.

\begin{figure}[htb]
    \centering
    \begin{tikzpicture}[x=1cm,y=1cm]
        \node[bloco, fill=corApp, minimum width=2.7cm] (prog) at (0,0) {Programa de\\teste (montador)};
        \node[bloco, fill=corCom, minimum width=2.7cm] (down) at (3.6,0) {Transferência\\pelo protocolo};
        \node[bloco, fill=corVM, minimum width=2.7cm] (run) at (7.2,0) {Execução na\\máquina virtual};
        \node[bloco, fill=corMem, minimum width=2.7cm] (obs) at (10.8,0) {Imagem de saídas\\(depuração)};
        \node[bloco, fill=corApp, minimum width=2.7cm] (est) at (0,-2.3) {Vetor de\\estímulos};
        \node[bloco, fill=corRos, minimum width=2.7cm] (orac) at (7.2,-2.3) {Oráculo (modelo,\\cálculo, norma)};
        \node[estado, minimum width=2.7cm, fill=corHW] (cmp) at (10.8,-2.3) {Comparação\\e veredito};
        \draw[seta] (prog) -- (down);
        \draw[seta] (down) -- (run);
        \draw[seta] (run) -- (obs);
        \draw[seta] (obs) -- (cmp);
        \draw[seta] (orac) -- (cmp);
        \draw[seta] (est) -- (orac);
        \draw[seta] (est.north) -- (prog.south);
        \draw[seta] (est.east) -- ++(0.6,0) |- ([yshift=-4mm]run.south west) -- (run.south west);
    \end{tikzpicture}
    \caption{Fluxo de execução e avaliação de um caso de teste}
    \label{fig:fluxo-validacao}

    \fonte{Elaborado pelo autor (2026).}
\end{figure}

O critério de aceitação depende do grupo: nas lógicas de contatos e bobinas exige-se igualdade bit a bit da imagem de saídas em todas as varreduras; nos temporizadores admite-se erro de até uma varredura mais a resolução da base de tempo; nas grandezas numéricas comparam-se os resultados com o cálculo do oráculo dentro da tolerância do tipo (exata para \textit{S32}, relativa para \textit{F32}).
Os Quadros~\ref{qua:testes-logica} a~\ref{qua:testes-integracao} organizam a cobertura pelo repertório de instruções do Quadro~\ref{qua:grupos-instrucoes}, de modo que cada grupo tenha ao menos um caso associado.
O conjunto completo é conduzido sobre o STM32F767ZI.
A coluna de situação distingue os casos já executados dos apenas especificados: os resultados dos primeiros constam do \ref{ape:pid} e do registro de ensaios do protocolo, e os demais passam a ser registrados nesses quadros à medida que são executados.

\begin{quadro}[htb]
    \centering
    \small\singlespacing\sloppy
    \begin{tabular}{|p{0.8cm}|p{3.2cm}|p{4.8cm}|p{2.8cm}|p{2.3cm}|}
        \hline
        ID & Cobertura & Procedimento & Oráculo & Situação \\
        \hline
        L1 & Contatos e lógica booleana: \texttt{LD}, \texttt{LDN}, \texttt{AND}, \texttt{ANDN}, \texttt{OR}, \texttt{ORN}, \texttt{XOR}, \texttt{XORN}, \texttt{NOT} & Tabela-verdade exaustiva com até três entradas & Tabela-verdade (modelo) & Especificado \\
        \hline
        L2 & Bobinas: \texttt{OUT}, \texttt{SET}, \texttt{RESET}, \texttt{TOGGLE} & Retenção do estado entre varreduras; alternância por borda & Modelo \textit{Stateflow} & Especificado \\
        \hline
        L3 & Bordas: \texttt{LDP}, \texttt{LDF}, \texttt{ANP}, \texttt{ANF}, \texttt{ORP}, \texttt{ORF} & Evento de duração de uma varredura; nível constante não gera evento & Diagrama de tempo & Especificado \\
        \hline
        L4 & Pilhas: \texttt{MPS}, \texttt{MRD}, \texttt{MPP}, \texttt{ANB}, \texttt{ORB} & Ramificações aninhadas até a profundidade máxima da pilha e além dela & Modelo \textit{Stateflow} & Especificado \\
        \hline
        L5 & Lógica de \textit{rung} composta & Selo, intertravamento e partida-parada, combinando L1 a L4 & Modelo \textit{Stateflow} & Especificado \\
        \hline
    \end{tabular}
    \caption{Casos de teste de lógica de contatos e bobinas}
    \label{qua:testes-logica}
    \fonte{Elaborado pelo autor (2026).}
\end{quadro}

\begin{quadro}[htb]
    \centering
    \small\singlespacing
    \begin{tabular}{|p{0.8cm}|p{3.2cm}|p{4.8cm}|p{2.8cm}|p{2.3cm}|}
        \hline
        ID & Cobertura & Procedimento & Oráculo & Situação \\
        \hline
        T1 & \textit{TON} & Tempo programado constante e por registrador; quatro bases de tempo; \textit{EN} interrompido antes do fim; alteração do tempo programado durante a contagem & Diagrama IEC 61131-3 & Especificado \\
        \hline
        T2 & \textit{TOF} & Idem; religamento durante a contagem & Diagrama IEC 61131-3 & Especificado \\
        \hline
        T3 & \textit{TP} & Pulso de duração fixa; reacionamento durante o pulso ignorado & Diagrama IEC 61131-3 & Especificado \\
        \hline
        C1 & \textit{CTU}, \textit{CTD}, \textit{CTUD} & Contagem por borda de subida; saídas nos limites; valor programado constante e por registrador & Norma IEC 61131-3 & Especificado \\
        \hline
        C2 & \textit{Reset} e carga (\textit{load}) & Prioridade entre \textit{reset}, carga e contagem; valor negativo & Norma IEC 61131-3 & Especificado \\
        \hline
    \end{tabular}
    \caption{Casos de teste de temporizadores e contadores}
    \label{qua:testes-temporizacao}
    \fonte{Elaborado pelo autor (2026).}
\end{quadro}

\begin{quadro}[htb]
    \centering
    \small\singlespacing
    \begin{tabular}{|p{0.8cm}|p{3.2cm}|p{4.8cm}|p{2.8cm}|p{2.3cm}|}
        \hline
        ID & Cobertura & Procedimento & Oráculo & Situação \\
        \hline
        N1 & Aritmética \textit{S32}: soma, subtração, multiplicação, divisão, módulo, mínimo, máximo, incremento, decremento, negação e valor absoluto & Valores nominais, zero e limites do inteiro de 32 \textit{bits} & Cálculo em \textit{MATLAB} & Especificado \\
        \hline
        N2 & Aritmética \textit{F32} (mesmas operações) & Valores nominais, muito pequenos e muito grandes; tolerância relativa & Cálculo em \textit{MATLAB} & Especificado \\
        \hline
        N3 & Casos-limite aritméticos & Divisão por zero (resultado nulo), módulo por zero (sem escrita), valor absoluto de \texttt{INT32\_MIN} (saturado), escalonamento com faixa de entrada nula (sem escrita) & Especificação da instrução (\ref{ape:isa}) & Especificado \\
        \hline
        N4 & Comparação \textit{S32} e \textit{F32}: igual, diferente, maior, menor, maior ou igual, menor ou igual & Fronteiras (igualdade, mais e menos uma unidade) & Cálculo em \textit{MATLAB} & Especificado \\
        \hline
        N5 & Conversão e escalonamento: \texttt{ITOF}, \texttt{FTOI}, \texttt{SCALE} & Pontos extremos e intermediários da faixa; operando constante e por registrador & Cálculo em \textit{MATLAB} & Especificado \\
        \hline
    \end{tabular}
    \caption{Casos de teste de blocos numéricos}
    \label{qua:testes-numericos}
    \fonte{Elaborado pelo autor (2026).}
\end{quadro}

\begin{quadro}[htb]
    \centering
    \small\singlespacing
    \begin{tabular}{|p{0.8cm}|p{3.2cm}|p{4.8cm}|p{2.8cm}|p{2.3cm}|}
        \hline
        ID & Cobertura & Procedimento & Oráculo & Situação \\
        \hline
        A1 & Entradas analógicas, saídas PWM e saída analógica & Erro em $N$ pontos da faixa & Instrumento de medição & Especificado \\
        \hline
        B1 & Bloco \textit{RAMP} & \texttt{EN} desligado; borda de subida com valor inicial; subida e descida com taxas distintas; saída no alvo & Modelo \textit{Stateflow} & Especificado \\
        \hline
        B2 & Bloco PID: sete cenários de unidade e regressão & Modos, transferência sem salto, banda morta, \textit{setpoint} dinâmico (\ref{ape:pid}) & Modelo de referência & Executado \\
        \hline
        B3 & Bloco PID em malha fechada (\textit{HIL}) & Cinco arquétipos de planta (\ref{ape:pid}) & Modelo da planta & Executado \\
        \hline
    \end{tabular}
    \caption{Casos de teste de entradas e saídas analógicas e de blocos de controle}
    \label{qua:testes-controle}
    \fonte{Elaborado pelo autor (2026).}
\end{quadro}

\begin{quadro}[htb]
    \centering
    \small\singlespacing
    \begin{tabular}{|p{0.8cm}|p{3.2cm}|p{4.8cm}|p{2.8cm}|p{2.3cm}|}
        \hline
        ID & Cobertura & Procedimento & Oráculo & Situação \\
        \hline
        I1 & Transferência de \textit{bytecode} & Programa recebido idêntico ao enviado; execução após a carga & Comparação byte a byte & Especificado \\
        \hline
        I2 & Robustez do protocolo & Comando inválido, tamanho inconsistente, CRC incorreto e identificador de dispositivo inválido, com o \textit{NACK} correto para cada categoria & Código de \textit{NACK} esperado & Executado \\
        \hline
        I3 & Robustez da máquina virtual & Opcode inválido leva a \textit{trap}; índice fora da faixa não acessa memória externa às regiões & Comportamento especificado & Especificado \\
        \hline
        P1 & Tempo de varredura & Programas de tamanho crescente; média, máximo e variação & Medição & Especificado \\
        \hline
        P2 & Consumo de memória & \textit{Flash} e RAM ocupadas, a partir do relatório de \textit{link} & Medição & Especificado \\
        \hline
        P3 & Custo por grupo de instruções & Tempo de execução por instrução em cada grupo & Medição & Especificado \\
        \hline
    \end{tabular}
    \caption{Casos de teste de integração, robustez e desempenho}
    \label{qua:testes-integracao}
    \fonte{Elaborado pelo autor (2026).}
\end{quadro}


%% ------------------------------------------------------------------------- %%
\section{Síntese}
\label{sec:sintese-implementacao}

Este capítulo apresentou a implementação da solução cujos fundamentos foram estabelecidos no Capítulo~\ref{cap:fundamentacao}: um motor de execução \textit{fetch-decode-execute} com palavra-base de instrução de 32 \textit{bits} e despacho em tempo constante; blocos funcionais com estado segregado em mapas de memória de acesso validado; uma camada de abstração que confina ao contrato de um único módulo toda a dependência do STM32F767ZI; um protocolo com enquadramento COBS e verificação CRC16; modelos de referência em \textit{Simulink}/\textit{Stateflow}; e uma estratégia de validação em quatro níveis, com cobertura de todos os grupos do repertório de instruções.
As escolhas de projeto e suas justificativas estão reunidas no Quadro~\ref{qua:escolhas-projeto}.

\begin{quadro}[htb]
    \centering
    \begin{tabular}{|p{4.7cm}|p{9.7cm}|}
        \hline
        Escolha & Justificativa \\
        \hline
        Palavra-base de 32 \textit{bits} com estrutura auxiliar opcional & Decodificação uniforme da primeira palavra \cite{patterson:2020}; constantes entram como palavras adicionais e parâmetros de blocos por instruções \texttt{ST} e \texttt{STW}, sem ampliar o formato base. \\
        \hline
        Despacho por tabela de opcodes & Tempo de despacho constante e independente da instrução \cite{liu:2000}. \\
        \hline
        Alocação estática de memória & Sem fragmentação nem variação de tempo de alocação \cite{liu:2000}. \\
        \hline
        \textit{Bytecode} somente leitura e estado em mapas & Permite atualizar a lógica sem regravar o \textit{firmware} e conter erros do usuário (\textit{sandbox}). \\
        \hline
        Validação de índice em cada acesso & Um erro de endereçamento não alcança memória fora das regiões alocadas. \\
        \hline
        \textit{Bytecode} validado e pré-processado na carga & Retira configuração estática do laço de varredura. \\
        \hline
        Aritmética \textit{S32} e \textit{F32} nativa & A FPU do alvo executa \textit{F32} em \textit{hardware} \cite{st:ds11532}. \\
        \hline
        Temporização pela base de tempo do sistema & O intervalo independe do tamanho do programa. \\
        \hline
        PID com amostragem fixa & A variação do tempo de varredura não entra nos termos integral e derivativo \cite{astrom:2006}. \\
        \hline
        Derivativo sobre a variável de processo & Evita o pico de saída quando o \textit{setpoint} muda \cite{astrom:2006}. \\
        \hline
        Anti-\textit{windup} por integração condicional & Suspende a integração sob saturação e a retoma na inversão do erro \cite{astrom:2006}. \\
        \hline
        Rastreamento do integrador no modo manual & Transição de manual para automático sem salto \cite{astrom:2006}. \\
        \hline
        Enquadramento COBS e CRC16 & Delimitação sem ambiguidade com sobrecarga limitada e detecção de corrupção \cite{cheshire:1999}. \\
        \hline
        Identificador de dispositivo nos quadros & Rejeição de comandos endereçados a outro dispositivo. \\
        \hline
        \textit{STM32Cube HAL} em vez de acesso direto a registradores & Reduz risco de erro de inicialização; sobrecusto contido na camada de abstração. \\
        \hline
        STM32F767ZI & FPU, memórias acopladas ao núcleo, periféricos e ferramental integrado \cite{st:ds11532,st:um1974}. \\
        \hline
    \end{tabular}

    \caption{Escolhas de projeto e justificativas}
    \label{qua:escolhas-projeto}
    \fonte{Elaborado pelo autor (2026).}
\end{quadro}
